\begin{pastexam}[label={px:ncu-108-dsa-07}]{中央資工所 108 學年度 第 7 題}
（複選；每題 5 分，答錯倒扣 1 分）

A hash function maps a key into a bucket in the hash table. Which of the following statements is true?
\begin{enumerate}[label=(\Alph*)]
\item A division hash function with divisor D = $7^r$, where $r$ is an integer, may result in serious collision.
\item Compared with hash chaining, open addressing requires fewer bucket accesses (in average)
\item When hash chaining is used to resolve overflows, the search for a key involves comparison with
keys that have different hash values.
\item In dynamic hashing, data buckets grows or shrinks (added or removed dynamically) as the records
increases or decreases.
\item None of the above.
\end{enumerate}
\end{pastexam}

\begin{pastsoln}
答案：(A)(D)。

\begin{itemize}
\item \textbf{(A) 對}。 除法 $h(k) = k \bmod D$ 的除數要避開小的質因數（資結書第~12.1~節）。 $D = 7^r$ 只有質因數
$7$；由資結書引理~12.1，鍵都是 $7$ 的倍數時，$k \bmod 7^r$ 也都是 $7$ 的倍數，只會用到 $D/7$ 個桶。 例如
$D = 49$，鍵 $7, 14, 21, \ldots$ 只落在 $0, 7, 14, \ldots, 42$ 這 $7$ 個桶。 實際資料常有這種規律（例如以 $7$ 為
週期的日期編號），所以會造成大量碰撞。 一般選不含小質因數的 $D$，最好是質數。
\item \textbf{(B) 錯}。 在同樣的載入因子 $\alpha$ 下，鏈結法平均看的比較少，不是比較多。 鏈結法失敗搜尋平均要
比較 $\alpha$ 次、成功搜尋約 $1 + \alpha/2$ 次（資結書第~12.2.2~節）；開放定址就算在最理想的均勻雜湊下，
失敗搜尋也要約 $1/(1 - \alpha) \ge 1 + \alpha$ 次，成功搜尋約
$\frac{1}{\alpha}\ln\frac{1}{1 - \alpha} = 1 + \frac{\alpha}{2} + \frac{\alpha^2}{3} + \cdots \ge 1 + \frac{\alpha}{2}$ 次（資結書
定理~12.5、12.7），線性探測更多。 開放定址的探測序列會經過別的鍵的家（home bucket），鏈結法只看自己那條串列。
\item \textbf{(C) 錯}。 鏈結法裡第 $i$ 條串列只放 $h(k) = i$ 的鍵，搜尋 $k$ 時只走 $h(k)$ 那一條，比較到的鍵雜湊值都
和 $k$ 相同。 會和雜湊值不同的鍵比較的是開放定址（例如線性探測走進別人的群聚裡）。
\item \textbf{(D) 對}。 動態雜湊（例如可延伸雜湊，資結書第~12.5~節）就是為了讓表跟著資料量變：桶滿了就拆成兩個，
資料變少時把相鄰的桶合併，不必整張表重建。
\item \textbf{(E)} (A)、(D) 對，所以錯。
\end{itemize}
\end{pastsoln}
